%--------------------------------------------------------------------
\subsection{Optical Response and Postprocessing}
\label{sec:optical-response}
%--------------------------------------------------------------------

The velocity-gauge formulation developed above is not restricted to an
impulsive perturbation.  For a prescribed spatially homogeneous vector
potential $\mathbf A(t)$, the formal Hamiltonian retains the external
field for all times at which the chosen electromagnetic perturbation
requires it,
\begin{equation}
  \hat H_V(t)
  =
  \frac{1}{2m_e}
  \left[
    \hat{\mathbf p}
    -
    \frac{q_e}{c}\mathbf A(t)
  \right]^2
  + \hat V_{\rm loc}[\rho](t)
  + \hat V_{\rm NL}^{V}(t).
  \label{eq:optical_response_vg_hamiltonian}
\end{equation}
Thus the same formal propagation can describe an impulsive excitation,
a finite-duration pulse, or a persistent time-dependent driving field,
provided the corresponding $\mathbf A(t)$ is used in the Hamiltonian and
in the current operator of Sec.~\ref{sec:current-nlpp}.

We use the Fourier-transform convention
\begin{equation}
  f(\omega)
  =
  \int_{-\infty}^{\infty}
  f(t)e^{-\imath\omega t}\,dt,
  \qquad
  f(t)
  =
  \frac{1}{2\pi}
  \int_{-\infty}^{\infty}
  f(\omega)e^{\imath\omega t}\,d\omega .
  \label{eq:fourier_convention}
\end{equation}
With this convention,
\begin{equation}
  \mathcal F\!\left[
    \frac{df}{dt}
  \right]
  =
  \imath\omega f(\omega),
  \label{eq:fourier_derivative_convention}
\end{equation}
assuming that boundary terms vanish or are treated separately.  In the
homogeneous velocity gauge,
\begin{equation}
  \mathbf E(t)
  =
  -\frac{1}{c}\frac{d\mathbf A(t)}{dt},
  \qquad
  \mathbf A(t)
  =
  \mathbf A(t_0)
  -
  c\int_{t_0}^{t}\mathbf E(t')\,dt',
  \label{eq:e_a_time_relation_optical}
\end{equation}
and therefore
\begin{equation}
  E_\alpha(\omega)
  =
  -\frac{\imath\omega}{c}
  A_\alpha(\omega).
  \label{eq:e_a_frequency_relation_optical}
\end{equation}

A conventional broadband electric-field kick used in real-time linear
response calculations\cite{bertsch-96,jakowski-rmg-tddft-2025} is
\begin{equation}
  E_\beta(t)
  =
  E_{0,\beta}\delta(t).
  \label{eq:delta_electric_field_optical}
\end{equation}
For $A_\beta(t<0)=0$, Eq.~\eqref{eq:e_a_time_relation_optical} gives
\begin{equation}
  A_\beta(t)
  =
  -cE_{0,\beta}\Theta(t).
  \label{eq:step_vector_potential_optical}
\end{equation}
Thus, in the formally equivalent velocity-gauge representation, the
vector potential remains at its post-kick value after the instantaneous
electric-field impulse.

The time-dependent observable is the physical macroscopic charge-current
density defined in Sec.~\ref{sec:current-nlpp}.  We write the induced
current as
\begin{equation}
  \Delta J_\alpha(t)
  =
  J_\alpha(t)-J_\alpha^{\rm eq},
  \label{eq:induced_current_definition}
\end{equation}
where $J_\alpha^{\rm eq}$ is the equilibrium current before the external
perturbation.  For the non-current-carrying equilibrium states considered
here, the macroscopic equilibrium current vanishes in the exact ground
state.  Retaining the induced-current notation is nevertheless useful in
numerical calculations, where finite $\mathbf k$-point sampling and
numerical errors may produce a small residual baseline.

In the linear-response regime,
\begin{equation}
  \Delta J_\alpha(\omega)
  =
  \sum_\beta
  \sigma_{\alpha\beta}(\omega)
  E_\beta(\omega),
  \label{eq:conductivity_linear_response}
\end{equation}
which defines the optical conductivity tensor
$\sigma_{\alpha\beta}(\omega)$.  Separate perturbations along linearly
independent Cartesian directions may be used to construct the full
tensor.  For the electric-field impulse of
Eq.~\eqref{eq:delta_electric_field_optical},
$E_\beta(\omega)=E_{0,\beta}$, and therefore
\begin{equation}
  \sigma_{\alpha\beta}(\omega)
  =
  \frac{\Delta J_\alpha^{(\beta)}(\omega)}
       {E_{0,\beta}},
  \label{eq:conductivity_from_delta_kick}
\end{equation}
where $\Delta J_\alpha^{(\beta)}$ denotes the current response to a kick
polarized along direction $\beta$.

The dielectric response follows from the relation between current and
macroscopic polarization.  Defining the induced polarization by
\begin{equation}
  \Delta J_\alpha(t)
  =
  \frac{d\Delta P_\alpha(t)}{dt},
  \label{eq:current_polarization_time}
\end{equation}
one obtains
\begin{equation}
  \Delta J_\alpha(\omega)
  =
  \imath\omega\Delta P_\alpha(\omega),
  \qquad
  \Delta P_\alpha(\omega)
  =
  \frac{\Delta J_\alpha(\omega)}{\imath\omega}.
  \label{eq:current_polarization_frequency}
\end{equation}
Thus integration of the current in time corresponds formally to division
by $\imath\omega$ in frequency space, with the zero-frequency component
requiring separate treatment.

Using Gaussian electromagnetic units, the displacement field satisfies
\begin{equation}
  D_\alpha(\omega)
  =
  E_\alpha(\omega)
  +
  4\pi P_\alpha(\omega)
  =
  \sum_\beta
  \varepsilon_{\alpha\beta}(\omega)E_\beta(\omega).
  \label{eq:gaussian_displacement_field}
\end{equation}
Combining Eqs.~\eqref{eq:conductivity_linear_response} and
\eqref{eq:current_polarization_frequency} gives
\begin{equation}
  \varepsilon_{\alpha\beta}(\omega)
  =
  \delta_{\alpha\beta}
  -
  \frac{4\pi\imath}{\omega}
  \sigma_{\alpha\beta}(\omega).
  \label{eq:dielectric_from_conductivity}
\end{equation}
With the Fourier convention of Eq.~\eqref{eq:fourier_convention}, the
dissipative diagonal response is represented by
${\rm Re}\,\sigma_{\alpha\alpha}(\omega)$ or, equivalently, by
$-{\rm Im}\,\varepsilon_{\alpha\alpha}(\omega)$.  The sign of the
imaginary part follows from the $e^{-\imath\omega t}$ transform
convention used here.

The conversion from dielectric response to an absorption coefficient or
absorbance depends on dimensionality and optical geometry.  For a
homogeneous bulk three-dimensional medium, one may introduce a complex
refractive index
$N_\alpha(\omega)=n_\alpha(\omega)+\imath\kappa_\alpha(\omega)$ along a
principal axis, with $N_\alpha^2=\varepsilon_{\alpha\alpha}$ and the
sign of $\kappa_\alpha$ chosen to describe attenuation of the physical
field.  The corresponding bulk absorption coefficient is
\begin{equation}
  \alpha_{{\rm abs},\alpha}^{\rm 3D}(\omega)
  =
  \frac{2\omega}{c}\kappa_\alpha(\omega).
  \label{eq:bulk_absorption_coefficient}
\end{equation}
Absorbance additionally depends on the optical path length and sample
geometry.  Low-dimensional periodic systems, such as monolayers treated
with supercell boundary conditions, require the appropriate
geometry-dependent conversion.  When only relative spectral features are
compared, normalized absorption-like spectra should therefore be
distinguished from quantitatively normalized conductivities, dielectric
functions, absorption coefficients, or absorbances.

In a finite real-time calculation, the Fourier transform is evaluated
over a finite propagation interval and is commonly combined with a
window or damping function to control ringing and define an effective
spectral broadening.  A generic finite-time estimate of the induced
current spectrum is
\begin{equation}
  \Delta J_\alpha^{T}(\omega)
  =
  \int_{0}^{T}
  W(t)\Delta J_\alpha(t)e^{-\imath\omega t}\,dt,
  \label{eq:finite_time_current_transform}
\end{equation}
where $W(t)$ is the chosen damping or window function.  The finite
propagation time and the choice of $W(t)$ determine the practical
frequency resolution and broadening of the computed spectrum.

For collinear spin-polarized calculations, the total charge current is
the sum of the spin-channel currents introduced in
Sec.~\ref{sec:collinear-spin-extension},
\begin{equation}
  \Delta J_\alpha(t)
  =
  \sum_{\sigma=\uparrow,\downarrow}
  \Delta J_\alpha^\sigma(t).
  \label{eq:spin_sum_induced_current_optical}
\end{equation}
The total optical conductivity follows from the total charge current in
Eq.~\eqref{eq:conductivity_linear_response}.  Spin-channel-resolved
spectra may also be formed from the separate
$\Delta J_\alpha^\sigma(t)$ responses within the fixed collinear spin
basis; this does not imply spin-flip transitions or noncollinear spin
dynamics.

Equations~\eqref{eq:e_a_time_relation_optical}--\eqref{eq:finite_time_current_transform}
define the formal current-based optical response used in the present
periodic velocity-gauge RT-TDDFT theory.\cite{Luber-rt-tddft}  The
numerical realization of the impulsive perturbation, the normalization
of raw current traces, and the finite-time processing used in RMG are
described in Sec.~\ref{sec:implementation}.
